% =====================================================================
% background.tex --- combined background draft for IWAI 2026 paper
% "Paradigm Dynamics in Multi-Agent Motivated Belief" (Hallgren & Hyland)
%
% Generated 2026-05-29 by combining 7 sub-agent drafts (S1_*.tex...S7_*.tex)
% in the same directory. Each \section corresponds to a section of the
% working-structure PDF (IWAI_26_Paradigm (1).pdf, 2026-05-29).
%
% Cross-references between sections have been normalised:
%   - sec:scalar-multiagent   (was sec:scalar in S5)
%   - sec:derived-results     (was sec:structured-results in S3)
%
% =====================================================================
% CONSOLIDATED KEY-HANDLING TODO (collected from all 7 drafts)
% =====================================================================
%
% A. Auto-keys to hand-rename before merging cluster bibs into refs.bib
% --------------------------------------------------------------------
% §1 + §4 (cluster_C6_theory_laden):
%   Chang2012_9         -> chang2012WaterChemRev                 % ANCHOR (Lavoisier/water)
%   Carrier2002_4       -> carrier2002IncommensPhlogiston
%   Malone1993_5        -> malone1993IncommensNoRelativism
%   Shimony1976_10      -> shimony1976KuhnTheses
%   Siegel1987_15       -> siegel1987KuhnInspired
%   Brewer2010_14       -> brewerLoschky2010TopDown
%   Lewowicz2011_2      -> lewowicz2011PurloinedReferent
%   Boantza2008_4       -> boantza2008PhlogisticHeat
%   Halloun2016_8       -> halloun2016ModelLadenInquiry
%   Mcevoy1989_7        -> mcevoy1989LavoisierLinguistic
%   Perrin1988_5        -> perrin1988ResearchTraditions
%
% §2 (cluster_C1_motivated_reasoning + cluster_C2_variational_cost):
%   Kahan2007_1         -> kahanEtAl2007WhiteMaleEffect          % ANCHOR
%   Kahan2017_4         -> kahan2017IdentityProtective           % ANCHOR
%   Kahan2017_2         -> kahanPetersDawson2017MotivatedNumeracy
%   Kahan2015_12        -> kahan2015ExpressiveRationality
%   Glüer-Pagin2024_11  -> gluerSpectre2024MotivationInNumeracy
%   Bolsen2019_7        -> bolsenPalm2019MotivatedReasoning
%   Little2020_13       -> littleSchnakenbergTurner2021MotivatedAccountability
%   Savolainen2022_3    -> savolainen2022HealthMisinformation
%   Kruglanski2020_10   -> kruglanski2020AllThinkingWishful
%   Hopkins2016_8       -> hopkins2016FreeEnergyDreaming
%   Friston2023_7       -> fristonEtAl2023FederatedInference
%   Costa2020_4         -> daCosta2020ActiveInferenceDiscrete
%
% §3 (cluster_C4_social_network_epistemology, intro subset):
%   Bala1998_1          -> balaGoyal1998LearningNeighbours        % ANCHOR
%   Bala2001_3          -> balaGoyal2001Conformism
%   Boudourides2026_7   -> boudourides2026FJBoundaryValue
%   Proskurnikov2015_6  -> proskurnikovMatveevCao2015Hostile
%   Ren2020_13          -> ren2020DeGrootCommunities
%
% §4 (cluster_C5_echo_chambers):
%   Nguyen2018_1        -> nguyen2020EchoChambers                  % ANCHOR
%   Pollock2024_3       -> pollock2024ContextualDiscordance
%   Clinkenbeard2025_4  -> clinkenbeard2025NihilisticDistrust
%
% §5 (cluster_C3_bmr_lakatos + cluster_C7_attention_precision):
%   Friston2018_2       -> fristonParrZeidman2018BMR              % ANCHOR
%   Thagard1989_1       -> thagard1989ExplanatoryCoherence        % ANCHOR
%   Thagard1988_4       -> thagardNowak1988ContinentalDrift
%   Ranney1988_3        -> ranneyThagard1988NaivePhysics
%   Markovi'c2023_1     -> markovicFristonKiebel2024BMRSparsify
%   Heijdra1986_6       -> heijdraLowenberg1986DuhemQuineLakatos
%   Wahyudi2026_2       -> wahyudi2026QuineWebOfBelief
%   Ostapiuk2025_11     -> ostapiuk2025LakatosResearchProgrammes
%   Mirza2019_2         -> mirzaAdamsFristonParr2019SelectiveAtt  % ANCHOR
%
% §6 (cluster_C4 derived-results subset + cluster_C5):
%   Griparic2021_7      -> griparicEtAl2021AlgebraicConnectivity
%   Montijano2017_11    -> montijano2017ConsensusSpectral
%   Tan2008_14          -> tan2008FiedlerOpinion
%   Müürsepp2019_7      -> muursepp2019FeyerabendPluralism
%   Novaes2025_1        -> novaes2025LonginoCritique
%   Ajdari2023_6        -> ajdari2023LonginoCritique
%   Borgerson2011_5     -> borgerson2011LonginoCritique
%   Longino1990_10      -> longino1990ScienceSocialKnowledge      % ANCHOR
%   Romero2017_2        -> romero2017NoveltyReplicability
%   Turner2023_2        -> turner2023ChamberTrust
%
% §7 (cluster_C8_algorithmic_curation):
%   D2025_2             -> tkach2025EthicsPersonalization
%   Lin2026_1           -> lin2026AlgorithmicBiasRecommender
%   Zhao2020_3          -> zhao2020InformationCocoons
%
% B. Foundational anchors MISSING from all clusters --- hand-add to refs.bib
% -------------------------------------------------------------------------
% §2:
%   kunda1990MotivatedReasoning           Kunda, Psych. Bull. 108(3), 1990
%   nickerson1998ConfirmationBias         Nickerson, Rev. Gen. Psych. 2(2), 1998
%   lordRossLepper1979BiasedAssimilation  Lord, Ross & Lepper, JPSP 37(11), 1979
% §5:
%   lakatos1970FalsificationMethodology   Lakatos, in Lakatos & Musgrave (eds.), 1970
%   quine1951TwoDogmas                    Quine, Phil. Review 60(1), 1951
%   fristonPenny2011PostHocBMS            Friston & Penny, NeuroImage 56(4), 2011
%   thagardVerbeurgt1998CoherenceConstraint  Thagard & Verbeurgt, Cog. Sci. 22(1), 1998
% §6:
%   zollman2007NetworkEpistemology        Zollman, Phil. Sci. 74(5), 2007
%   zollman2010TransientDiversity         Zollman, Erkenntnis 72, 2010
%   kitcher1990DivisionCognitiveLabor     Kitcher, J. Phil. 87(1), 1990
%   strevens2003PriorityRule              Strevens, J. Phil. 100(2), 2003
%   muldoonWeisberg2011Robustness         Muldoon & Weisberg, Phil. Stud. 152(2), 2011
%
% =====================================================================


% =====================================================================
\section{Theory-Laden Perception}
\label{sec:theoryladen}
% =====================================================================

\paragraph{Kuhn's three faces.}
The central phenomenological commitment of this paper is the
Kuhnian observation that the paradigm a scientist holds shapes what
she is even capable of perceiving \citep{kuhn1962Structure}. This is
not a methodological lapse to be corrected by more careful
experimental design but a structural feature of scientific cognition,
and it has three interconnected faces that we will treat as
expressions of a single underlying phenomenon. The first face is
what a scientist sees in incoming data: the same observation is read
differently from inside different paradigms, with the assignment of
\emph{anomaly} versus \emph{nuisance} itself being paradigm-dependent
\citep{Brewer2010_14, Shimony1976_10}. The second face is what
experiments she chooses to run: the paradigm shapes the space of
investigations she finds worth undertaking, which in turn shapes
what evidence the community as a whole produces
\citep{Halloun2016_8}. The third face is what theoretical
alternatives she even considers: rival frameworks are not so much
rejected as never seriously entertained, with the candidate set of
live revisions itself being structured by the prevailing framework
\citep{Carrier2002_4, Malone1993_5}.

\paragraph{Theory-ladenness as structural feature, not bias.}
The contemporary literature has substantially clarified what was, in
Kuhn's original framing, sometimes left ambiguous. \citet{Siegel1987_15}
and \citet{Malone1993_5} reconstruct incommensurability as a
condition on referential and inferential overlap between frameworks
rather than as a relativist denial of cross-paradigm rationality, and
\citet{Carrier2002_4} demonstrates on the phlogiston case that
\emph{empirical comparability} between rival paradigms is possible
without their being intertranslatable term-by-term. Cognitive
psychology has independently corroborated the perceptual layer of the
claim: \citet{Brewer2010_14} review top-down and bottom-up
contributions to scientific observation and find that the prior model
substantively conditions what counts as a datum, in line with the
phenomenology Kuhn was describing. The resulting picture is
deflationary in one direction --- theory-ladenness is not the
collapse of objectivity --- and inflationary in another: the
influence of the prevailing framework reaches further into the
cognitive pipeline than an additive bias term would suggest, touching
attention, salience, the felt space of alternatives, and the
allocation of experimental effort.

\paragraph{The Chemical Revolution as a worked case.}
The Chemical Revolution is the canonical historical demonstration of
all three faces operating jointly. \citet{Chang2012_9} argues that
phlogiston theory was not a failed predecessor refuted by Lavoisier's
data but a research programme with its own coherent set of
observables, anomalies, and live revisions; its replacement was less
a victory of evidence over confusion than a reallocation of what
counted as a phenomenon to be explained.
\citet{Lewowicz2011_2} and \citet{Boantza2008_4} reinforce this
reading from inside the chemistry: the phlogiston framework had
resources --- including a treatment of heat
\citep{Boantza2008_4} --- whose referents were not so much falsified
as relocated under the oxygen framework, exactly the kind of move
that the third Kuhnian face predicts. This historical record matters
for our model because it shows that the three faces are not
independent variables to be added: shifts in what is observed, what
is investigated, and what is entertained co-vary, and they co-vary
on the timescale of community-wide reorganisation rather than
individual reasoning.

\paragraph{Methodological commitment.}
We take this as the empirical commitment that licenses the model's
complexity. An accumulated-checklist alternative --- enumerate
biases in perception, biases in experiment choice, and biases in
hypothesis generation, then sum them --- would be both more
defensive and less informative, because it would treat the three
faces as separable failure modes rather than as expressions of a
single structural fact about how a community-held framework
conditions cognition. The stronger rhetorical position, and the one
this paper takes, is that a single underlying mechanism is
responsible for all three faces, and that the model must therefore
couple perception, action, and hypothesis-space in a way that an
additive bias account cannot. The remainder of the paper develops
that mechanism. In our setting, the
shared framework is encoded in the joint structure of agents'
generative models and the channel of inter-agent precision that
binds them; the three faces will reappear, respectively, as the
likelihood-weighting in inference, as the experiment-selection
policy, and as the support of the prior over theories. The case
for treating them together, rather than enumerating them
separately, rests on the phenomenological commitment defended
here.


% =====================================================================
\section{Motivated Belief at the Single-Agent Level}
\label{sec:motivatedbelief}
% =====================================================================

\paragraph{Empirical motivated reasoning.}
The cognitive-psychological case that belief revision is shaped by
goals other than accuracy is long-standing and converged. The
canonical statement is Kunda's argument that reasoning is steered by
\emph{directional} as well as accuracy goals: subjects construct the
sequence of inferential steps that licenses the conclusion they would
prefer to reach, subject to the constraint that the construction
remain plausible to themselves \citep{kunda1990MotivatedReasoning}.
This framing absorbed two earlier empirical findings as special
cases: the biased-assimilation experiments of
\citet{lordRossLepper1979BiasedAssimilation}, in which the same
mixed evidence about a contested policy strengthened the prior
commitments of partisans on \emph{both} sides; and confirmation-bias
phenomena more broadly catalogued by \citet{nickerson1998ConfirmationBias}.
The empirical literature has since accumulated in two directions.
The first is Kahan's cultural-cognition programme, in which the
directional goal is identified specifically as \emph{identity
protection}: subjects process risk-relevant evidence in whatever way
preserves their standing in the affinity groups they belong to
\citep{Kahan2007_1, Kahan2017_4}. The signature empirical finding is
that the gap between cultural sub-populations \emph{widens} with
numerical and scientific sophistication rather than narrowing
\citep{Kahan2017_2}; Kahan reads this as evidence that the more
capable reasoners are more effective at the selective evaluation of
identity-incongruent evidence, a behaviour he characterises as
\emph{expressive rationality} --- rational relative to the agent's
practical stakes in group membership, irrational relative to truth
\citep{Kahan2015_12}. The mechanism remains contested:
\citet{Glüer-Pagin2024_11} argue that what Kahan's experiments
isolate is motivation \emph{to reason} rather than motivated reasoning
in the strict sense, and the broader political-science literature now
distinguishes carefully between directional and accuracy regimes and
the conditions under which one or the other dominates
\citep{Bolsen2019_7, Little2020_13, Savolainen2022_3}. For our
purposes the contested points do not matter: the robust observation
is that human belief revision exhibits a systematic, content-sensitive
asymmetry in its receptivity to evidence, and that this asymmetry is
not adequately captured by treating the agent as an unbiased Bayesian
updater.

\paragraph{Variational formalisation.}
The active-inference community has begun to give this asymmetry a
principled mathematical form. The starting point is the variational
free-energy functional, in which an agent's update from prior to
posterior incurs a complexity cost proportional to the
Kullback--Leibler divergence between the candidate posterior and the
prior \citep{Costa2020_4}. \citet{hylandAlbarracin2025MindChange}
re-foreground this complexity term as a \emph{motivational} quantity:
the cost of changing one's mind is the KL divergence from the current
belief to the candidate revision, weighted by an affective utility
that reflects how much the agent cares about retaining the prior.
Because that divergence is asymmetric, belief revision is
path-dependent: evidence that would update a neutral prior into a
particular posterior need not, played in reverse, undo a revision
already made. This is the formal counterpart of the
biased-assimilation and identity-protection findings: preferred
beliefs are not merely held more confidently, they are protected by a
mathematically explicit revision cost that contrary evidence must
overcome. \citet{Kruglanski2020_10} reaches a sympathetic conclusion
from the social-psychological side, arguing that there is no sharp
boundary between motivated and non-motivated cognition because
\emph{all} inference is directional in the sense that it is steered
by what the agent values; the variational reading makes this claim
precise by locating the directional weight in the affective utility
that multiplies the complexity term. Adjacent strands of the
active-inference literature provide the surrounding machinery on
which the multi-agent extension will draw, including federated
inference under shared world models \citep{Friston2023_7} and the
free-energy account of belief-revision dynamics more generally
\citep{Hopkins2016_8}.

\paragraph{Open question.}
The single-agent picture is now reasonably well developed, both
empirically and formally. What it does not yet capture is the
phenomenon Kuhn observed at the community scale: that scientific
sub-fields exhibit collective motivated belief, sustained across
generations, with the directional goal residing not in any individual
ego but in the institutional and social structures that bind the
community together. The active-inference treatment so far has been
single-agent --- the cost of mind-change is borne by one mind, against
its own prior. The empirical motivated-reasoning literature is in
practice multi-agent --- the identity being protected is a group
identity --- but its formalisation has remained at the level of
within-individual mechanism. The question this paper takes up is how
the variational cost of belief revision generalises to a population
of coupled inferential agents in a way that recovers what Kuhn
observed: paradigm-level persistence that is sociologically generated,
asymmetric in time, and not reducible to the directional preferences
of any single agent.


% =====================================================================
\section{From Single-Agent to Multi-Agent: The Scalar Treatment}
\label{sec:scalar-multiagent}
% =====================================================================

\paragraph{Scalar Bayesian updating on a graph.}
The scalar multi-agent setting we adopt sits inside a well-developed
tradition that treats opinion change as a linear, convex-combination
update on a weighted directed graph. The canonical primitive is the
DeGroot averaging model \citep{Boudourides2026_7}, in which each
agent's opinion at the next step is a row-stochastic mixture of its
neighbours' current opinions; under mild connectivity assumptions this
contracts to consensus. The Friedkin--Johnsen extension breaks that
asymptotic consensus by allowing each agent to retain weight on its
initial belief through a susceptibility coefficient
$s_i \in [0,1]$, so that the long-run opinion vector is an
affine-resolvent functional of the boundary of stubborn agents
\citep{Boudourides2026_7}. The conservatism parameter in our setup
plays the same formal role as the Friedkin--Johnsen susceptibility,
read off the variational free-energy decomposition rather than
postulated: the cost-of-mind-change term in the tempered-Bayesian
update of \citet{hylandAlbarracin2025MindChange} is precisely an
asymmetric KL penalty on revising one's posterior, and at the linear
limit its scalar weight maps to $1 - s_i$. The literature has worked
out a great deal of structural detail on this family: closed-form
steady states via resolvent / Green-function representations,
quantitative convergence rates controlled by the spectral radius of
the susceptibility-weighted influence matrix, and explicit sensitivity
of the steady state to perturbations of either the influence matrix or
the susceptibility profile \citep{Boudourides2026_7}. Signed
generalisations recover polarisation as a stable attractor when the
underlying graph admits a structural-balance decomposition into
hostile camps \citep{Proskurnikov2015_6}, and community-detection
methods exploit the same consensus dynamics to read off hierarchical
group structure from opinion trajectories \citep{Ren2020_13}. We
inherit this machinery and re-read it in active-inference terms: the
influence matrix is a learned channel precision rather than a fixed
social-tie weight, and the susceptibility / conservatism is a
free-energy quantity rather than a phenomenological parameter.

\paragraph{Multi-armed bandit framing of theory choice.}
The complementary tradition treats the agent's problem not as
averaging an opinion but as choosing which arm of a multi-armed
bandit to pull next given its neighbours' payoffs, which we draw on
explicitly when the scalar opinion is read as a paradigm \emph{choice}
rather than a continuous belief. \citet{Bala1998_1} introduce the
canonical model: each agent repeatedly chooses among a finite set of
risky actions, observes its own and its neighbours' payoffs, and
updates a Bayesian posterior over which action is best. They prove
that on connected networks the population converges to consensus on a
single action almost surely, but the consensus action need not be the
optimal one --- early random fluctuations get locked in by social
imitation, generating a long tail of pathological lock-ins on
inferior arms. Their follow-up \citep{Bala2001_3} sharpens this into
a quantitative tradeoff between conformism and diversity that
prefigures the network-epistemology framing later picked up by
\citet{zollman2007NetworkEpistemology}, whose paradoxical-stubbornness and
communication-structure results we treat as derived consequences of
the same primitives and develop in
Section~\ref{sec:derived-results}. The Bala--Goyal lineage is the
right reference frame for our §3 simulations because it is the
minimal model in which (i) the agent has a well-defined Bayesian
posterior over a finite paradigm set, (ii) communication enters
through observed payoffs rather than through opinion averaging, and
(iii) the timescale-separation between learning and convergence
generates path-dependent population outcomes --- the three features
the scalar paradigm experiments are designed to exhibit.

\paragraph{Active-inference coupling.}
Bringing the two strands together in an active-inference register is
the contribution of \citet{albarracin2022EpistemicCommunities}, who
couple variational-Bayesian agents through shared observations and
characterise echo-chamber bifurcation as a function of the
inter-agent precision. Their construction is the immediate ancestor
of our §3 setup, with two refinements. First, the inter-agent
precision is itself learned, not fixed, so the trust matrix is the
only Bayesian-learned relational object in the model and the
``susceptibility profile'' of the Friedkin--Johnsen analysis becomes
endogenous to the dynamics. Second, the tempered-Bayesian update of
\citet{hylandAlbarracin2025MindChange} contributes an asymmetric
cost-of-mind-change term that does not appear in the symmetric
linear-averaging tradition but that is necessary to recover Kuhnian
path-dependence at the population level: even a homogeneous network
with isotropic trust will, with non-zero conservatism, lock onto its
initial-condition basin in a way that pure DeGroot or Bala--Goyal
dynamics cannot. The scalar experiments that follow probe how this
conservatism parameter interacts with the trust-precision topology
and with the affective utility on paradigm values to produce the
population behaviours we will then re-read, in §4, as the structural
limits of a scalar treatment.


% =====================================================================
\section{What the Scalar Treatment Cannot Capture}
\label{sec:scalar-limits}
% =====================================================================

\paragraph{What the scalar treatment recovers.}
Before locating its limits, it is worth being explicit about what the
scalar treatment of paradigm commitment already captures. A
temperature-modulated update on a single belief variable produces an
asymmetric assimilation: the same datum is integrated by an agent who
sits near one attractor with a different effective Bayes factor than by
an agent who sits near another, so that observers in different
paradigms reach different posteriors from a shared
observation~\citep{kuhn1962Structure}. This is not a degenerate
re-description of confirmation bias; it is a principled consequence of
the cost of mind-change, and it does some of the work that Kuhn's first
face of theory-ladenness demands. To that extent, the scalar treatment
is doing real explanatory labour, and we resist the temptation to throw
it away in order to motivate what follows. The point of this section is
to mark where its explanatory reach ends.

\paragraph{Structural reorganisation: the Lavoisian case.}
The clearest pressure on a scalar treatment comes from the Chemical
Revolution. \citet{Chang2012_9} argues that the disagreement between
Lavoisier and the late phlogistonians cannot be cashed out as a
disagreement over the probability assigned to a shared proposition.
What counted as combustion, what counted as a substance, what
counted as the conservation of something, and what an experimental
yield was even evidence \emph{for}, all shifted together; the very
referents of the disputants' terms drift across the
transition~\citep{Lewowicz2011_2}. \citet{Carrier2002_4} makes the
related point that intertheoretic comparison across the phlogiston/
oxygen divide requires reconstructing each side's network of
inferential commitments rather than aligning a common variable; the
incommensurability at stake is between generative models, not between
parameter settings within one. Historical work on the linguistic and
institutional reorganisation of Lavoisian
chemistry~\citep{Mcevoy1989_7,Perrin1988_5,Boantza2008_4} reinforces
this: the revolution was a re-tooling of which questions could be
posed, not a re-weighting of answers to fixed questions. A scalar
order parameter cannot, by construction, carry this information,
because the cardinality of admissible nodes and edges is exactly what
is changing. Even philosophers sympathetic to a deflationary reading
of incommensurability concede that the residue which resists
deflation is structural rather than
quantitative~\citep{Malone1993_5}. The scalar treatment is therefore
silent on what the historians take to be the substantive content of
the revolution.

\paragraph{Failure modes the scalar treatment cannot diagnose.}
A second pressure comes from a different direction. Even if one set
aside the question of structural reorganisation and asked only after
the social pathologies a multi-agent paradigm model ought to recover,
the scalar treatment runs into a distinction it cannot draw.
\citet{Nguyen2018_1} argues that two pathologies of collective
inference -- routinely conflated in popular usage -- have entirely
different mechanisms and admit entirely different remedies. An
\emph{epistemic bubble} is a failure of \emph{informational
topology}: relevant voices are omitted from the network, coverage is
incomplete, and exposure to those voices is sufficient to dissolve
the pathology. An \emph{echo chamber} is a failure of \emph{trust
calibration}: outside voices are not absent but actively discredited,
so that exposure to them, if anything, reinforces the chamber. The
two failure modes can co-occur, but they are not the same object;
follow-up work emphasises that they decouple in
practice~\citep{Pollock2024_3} and that the trust-calibration variant
can become self-stabilising in ways that omission alone cannot
explain~\citep{Clinkenbeard2025_4}. A scalar belief variable updated
under socially-weighted observations registers both pathologies as
the same kind of thing -- a population whose marginal has collapsed
onto one attractor -- and so cannot, on its own resources,
distinguish a community that has merely failed to encounter an
alternative from one that has built the discrediting of alternatives
into its updating rule. The two patients present with the same
symptom, but they have different diseases, and the scalar model has
no diagnostic on which to separate them.

\paragraph{Locating the gap.}
The Lavoisian case and the bubble/chamber distinction press from
opposite ends on the same point. From the philosophy-of-science side,
the object being reorganised in a paradigm shift is the inferential
web, and a scalar carries none of the relevant structure. From the
social-epistemology side, the pathologies a multi-agent treatment
ought to distinguish are differentiated by their structural
locus -- topology versus trust -- and a scalar cannot mark which
locus is failing. These are not arguments against the scalar
treatment as such; they are arguments that something has to be added
to it. What that something is, and how it can be added without giving
up the dynamical traction the scalar treatment already buys, is the
business of the next section.


% =====================================================================
\section{Structural Extension: Bayes Nets and Bayesian Model Reduction}
\label{sec:structural}
% =====================================================================

The scalar treatment of \S\ref{sec:scalar-multiagent} sufficed to expose
asymmetric belief-revision dynamics, but it cannot speak to what
practising scientists actually argue about: which commitments are
foundational, which are negotiable, and which channel of observation
to spend a finite attentional budget on next. The structured extension
recovers all three by promoting each agent's paradigm from a point in
$\Theta$ to a Bayesian network whose nodes carry theoretical
commitments and whose edges encode inferential dependencies.

\paragraph{Structural belief representation.}
The conceptual ancestry of paradigms-as-networks predates active
inference by half a century. \citet{Wahyudi2026_2}, reconstructing
Quine's coherentism, summarises the web-of-belief picture in which no
proposition confronts experience in isolation: revisions propagate
along inferential links, and the choice of which link to cut is
underdetermined by data alone. Lakatos's methodology of scientific
research programmes sharpens this into a topology with a designated
\emph{hard core} of commitments shielded by a flexible
\emph{protective belt} of auxiliary hypotheses
\citep{Heijdra1986_6,Ostapiuk2025_11}; the belt absorbs anomalies so
that the core remains in force, and a programme degenerates only when
belt-patching ceases to predict novel facts. Thagard's ECHO programme
operationalises Quinean coherence as a connectionist constraint
network and demonstrates, on the continental-drift case study, that
the same machinery can reconstruct a historical paradigm shift as a
re-equilibration of explanatory links
\citep{Thagard1989_1,Thagard1988_4}; \citet{Ranney1988_3} extend the
account to belief revision in naive physics. None of these accounts
furnishes a generative model, but each picks out the same structural
object: a graph of commitments whose densely-connected interior is
held rigid while its periphery negotiates with observation. We adopt
this object verbatim, identifying core nodes with the parameters
under tight prior precision and belt nodes with those under broad
prior precision in the agent's generative model.

\paragraph{Bayesian model reduction.}
Promoting the core/belt distinction to a piece of formal machinery
requires a principled treatment of \emph{structural} inference --- the
problem of scoring alternative graphs, not just alternative parameter
values, against the same data. \citet{Friston2018_2}, extending the
post-hoc model selection scheme of Friston \& Penny (2011), provide
precisely this in the form of Bayesian model reduction (BMR). The
construction is straightforward in variational Bayes: once the full
model has been inverted, the evidence and posterior of any
\emph{reduced} model --- one differing only in its priors --- can be
computed analytically from the full posterior and the ratio of priors,
generalising the Savage--Dickey density ratio. Concretely, switching
off a candidate edge corresponds to collapsing the prior variance on
the associated parameter to zero; the resulting change in variational
free energy is the (log) Bayes factor for that structural revision.
For the Dirichlet posteriors that govern categorical likelihood
mappings in discrete state-space models --- the natural form for
theory--observation links --- the reduced free energy has a closed
form in terms of multivariate beta functions. \citet{Friston2018_2}
showcase BMR as both a data-analytic tool and a metaphor for
neurobiological abductive reasoning: the same equations that prune
redundant DCM connections describe the synaptic elimination an agent
performs when ``thinking off-line'' simplifies its model of a
contingency. \citet{Markovi'c2023_1} extend the scheme to deep-network
sparsification, evidencing that the reduction is general across
parameter spaces. We invoke BMR as the engine of paradigm dynamics:
candidate revisions to an agent's network --- adding a belt link,
demoting a core node, accepting a rival auxiliary --- are scored by
their reduced free energy, with the structural affective utility
inherited from \S\ref{sec:motivatedbelief} entering as a contribution
to the cost of admitting each revision. Path-dependence in the scalar
model thus refines into something graph-theoretic: revisions that
touch the core are more expensive because more of the posterior has
to be re-balanced.

\paragraph{Active sampling and precision allocation.}
The structural representation closes the loop with action selection in
a way the scalar model could not. Observation in the structured model
is not exogenous arrival of evidence but an active sampling decision:
the agent allocates a finite precision budget across belt nodes,
sampling some channels with high precision and others with low.
\citet{friston2015ActiveInferenceEpistemic} establishes the
underlying primitive --- expected free energy decomposes into
extrinsic (utility) and epistemic (information-gain) value ---
and \citet{millidge2021EFE} surveys the equivalent decompositions.
\citet{Mirza2019_2} show, in an MDP formulation of Yarbus's classic
visual-search paradigm, that selective attention is just
context-dependent modulation of the likelihood precisions $\zeta$:
raising $\zeta$ on task-relevant channels and lowering it on
task-irrelevant ones reweights expected information gain so that the
agent saccades to evidence that distinguishes hypotheses it actually
entertains. The same construction underwrites
\citet{feldmanFriston2010Attention}, \citet{parrFriston2017Attention},
and the precision-control account of
\citet{parrFriston2019AdvancedAtt}; \citet{limanowskiFriston2018Dark}
fills in the failure mode where mis-allocated precision dampens
exploration. Transposed to the paradigm setting, precision allocation
\emph{is} the paradigm-dependence of observation: an agent inhabiting
a Newtonian network attends to mass-and-force belt nodes with high
precision and assigns negligible precision to channels whose only
hidden cause lies in a rival network. The result is a structured
analogue of theory-laden perception
(\S\ref{sec:theoryladen})~\citep{hylandAlbarracin2025MindChange},
now derived rather than stipulated: what counts as a salient
observation is what the agent's current network has bandwidth to
believe.


% =====================================================================
\section{Derived Consequences of the Structured Model}
\label{sec:derived-results}
% =====================================================================

The structural extension developed in Section~\ref{sec:structural} is
not introduced for its own sake. It earns its place because three
otherwise-disparate results in network epistemology --- the Zollman
tradeoff between consensus speed and asymptotic accuracy, the
paradoxical-stubbornness finding from the division-of-cognitive-labor
literature, and Nguyen's distinction between epistemic bubbles and
echo chambers --- recover as derived consequences of one mechanism: a
trust-precision network propagating candidate \emph{structural}
revisions rather than scalar opinions. We position each result in turn
and indicate which observable in our simulations it indexes.

\paragraph{The Zollman tradeoff and the value of inefficiency.}
\citet{zollman2007NetworkEpistemology, zollman2010TransientDiversity}
showed, in a multi-armed-bandit network of Bayesian scientists, that
\emph{denser} communication graphs converge faster but more often
lock onto the worse arm; sparser graphs sustain transient diversity
long enough for the inferior hypothesis to be discarded, and so reach
higher asymptotic accuracy. The mechanism is now well understood as a
spectral one: the rate at which beliefs synchronise on a graph is set
by its algebraic connectivity (the Fiedler value $\lambda_2$ of the
graph Laplacian)
\citep{Griparic2021_7, Montijano2017_11, Tan2008_14}, so larger
$\lambda_2$ compresses the transient phase during which dissent can
falsify a leading hypothesis. The structured model inherits this
argument unchanged at the topological level, but reroutes its
substantive content: the channels carry \emph{precision-weighted
candidate revisions} of the generative model, not scalar credences,
so the speed--accuracy tradeoff applies to structural revision events
themselves. This vindicates the Feyerabendian-pluralist intuition
\citep{Müürsepp2019_7} in mechanistic form --- a measured inefficiency
in the trust-precision network is the substrate condition for the
critical-discursive uptake that
\citet{longino1990ScienceSocialKnowledge} requires of an epistemically
healthy community --- and supplies the formal claim Longino's
critical contextual empiricism has been criticised for leaving
underspecified \citep{Novaes2025_1, Ajdari2023_6, Borgerson2011_5}.

\paragraph{Division of cognitive labour and paradoxical stubbornness.}
\citet{kitcher1990DivisionCognitiveLabor} and
\citet{strevens2003PriorityRule} argued that the priority rule
incentivises credit-driven scientists to spread across research
programmes, producing a community-level distribution of effort closer
to the social optimum than truth-driven scientists would individually
choose. \citet{muldoonWeisberg2011Robustness} relaxed the
god's-eye information assumption: in their local-information variant,
agents that appear \emph{paradoxically stubborn} --- unwilling to
abandon a low-success-probability programme --- in fact regularise
the system, preventing premature convergence on whichever programme
happened to lead first. \citet{romero2017NoveltyReplicability}
qualifies the optimistic Kitcher--Strevens reading by noting that the
priority rule rewards novelty regardless of replicability and so
\emph{disincentivises} the very direct-replication work that would
discipline structural revisions. In the structured model these
findings appear as a single statement: an agent's resistance to a
structural revision is proportional to the precision of its current
generative model, so agents with strong structural utility on their
incumbent paradigm act as protective elements against premature
structural collapse. They are not uniquely irrational; their
structural commitments simply require correspondingly stronger
cross-channel evidence to revise. The Muldoon--Weisberg paradox and
the Romero replication-vice fall out as two readings of the same
asymmetry --- one casting it as virtue (transient diversity
preserved), the other as vice (replication suppressed) --- and the
model exposes the parameter that interpolates between them.

\paragraph{Bubbles versus chambers: failure modes of cross-agent channels.}
\citet{nguyen2020EchoChambers} draws a distinction that, in our
framework, is structural rather than rhetorical. An
\emph{epistemic bubble} is a failure of \emph{trust topology}:
relevant voices are omitted, so the cross-agent channel is sparse or
disconnected, and the bubble pops on simple exposure. An
\emph{echo chamber} is a failure of \emph{trust precision}: the
discrediting beliefs assigned to outsiders pre-emptively downweight
incoming cross-surprisals, so the chamber survives exposure and can
even use it for reinforcement
\citep{Pollock2024_3, Turner2023_2, Clinkenbeard2025_4}. In the
scalar-opinion setting this distinction must be argued for
philosophically; in the structured-belief setting it is recovered as
the formal separation between sparsity of $\trustmat$ and the
precision weights it carries on its non-zero entries. Crucially,
because the channels now carry candidate \emph{structural} revisions
rather than scalar opinions, the bubble--chamber distinction becomes
more honest: a bubble withholds candidate revisions; a chamber
admits them and then assigns them zero precision. The two failure
modes call for different repairs --- coverage versus
precision-recalibration --- and the model makes the repair
prescription mechanistic rather than virtue-theoretic.


% =====================================================================
\section{Discussion}
\label{sec:discussion-background}
% =====================================================================
% Background prose for §7 of the working-structure PDF. Currently only
% the third open question (external attentional curation) is drafted;
% other open-question paragraphs are still authors' work.

\paragraph{External attentional curation.}
The third open question is the most consequential for the model's empirical
domain of validity, and it remains open. Throughout
Sections~\ref{sec:structural}--\ref{sec:derived-results} an agent's policy
over which neighbours, papers, and observations to attend to is treated as
endogenous: it is selected by expected free energy under the agent's own
generative model and affective utility. In contemporary scientific practice
this assumption is already strained. Citation recommenders, search-engine
ranking, social-media amplification, and journal-level visibility filters
increasingly mediate which findings an individual scientist ever encounters,
and the same logic operates a fortiori in adjacent epistemic communities ---
policy analysis, public discourse around scientific results, the training
corpora of large models that themselves now participate in summarisation and
triage. The literature on algorithmic curation documents that such systems
do not merely accelerate access to a fixed information environment but
actively reshape exposure through popularity bias, preference reinforcement,
and the formation of ``information cocoons'' that insulate users from
disconfirming evidence \citep{Lin2026_1, Zhao2020_3}, and that the boundary
between helpful personalisation and manipulation is structurally hard to
draw once the recommender's objective is decoupled from the user's epistemic
interest \citep{D2025_2}.

For the present construction the consequence is precise but unresolved. The
attentional policy $\pi$ no longer factorises through a single agent's
generative model: a curator $\mathcal{C}$ --- algorithmic, institutional, or
both --- intervenes between the latent paradigm field and the per-agent
observation channel, with its own objective that is in general not aligned
with predictive accuracy on the underlying scientific question. The
trust matrix $\trustmat$ then learns precisions over a doubly-mediated
signal, and the resource-flow readouts acquire a third upstream lever beyond
the exogenous inflow rate and initial trust configuration:
the curator's policy itself. Whether this collapses to a renormalisation of
the existing trust dynamics, or whether it introduces qualitatively new
attractors --- in particular, attractors that survive the corrective force
of cross-surprisal because the disconfirming observations are never sampled ---
is the model's natural next question. It is also the point at which the
construction makes contact with broader concerns about AI systems gradually
displacing human epistemic agency \citep{kulveit2025GradualDisempowerment} and
about the multi-principal coordination problems that arise when curators,
scientists, and institutions optimise different objectives over a shared
information substrate \citep{critchKrueger2020ARCHES}. We flag this as the
open direction the framework most naturally extends into, and the one we
take to be most urgent.
